% Part: second-order-logic
% Chapter: syntax-and-semantics
% Section: expressive-power

\documentclass[../../../include/open-logic-section]{subfiles}

\begin{document}

% SEGMENT OLP-0328-S01

\olfileid{sol}{syn}{exp}
\olsection{வெளிப்படுத்தும் ஆற்றல்}

\begin{explain}
இரண்டாம்தர மாறிகள்மீதான அளவையிடல், முதல்தரத் தருக்கத்தின்
வெளிப்படுத்தும் ஆற்றலைவிட மொழியின் ஆற்றலைப் பெருமளவு கூட்டுகிறது.
சில பண்புகளைக் கொண்ட சார்புகள் அல்லது உறவுகள் உள்ளன என்று
இரண்டாம்தர இருத்தலளவையிடல் கூற உதவுகிறது. முதல்தரத் தருக்கத்தில்
இதைச் செய்வதற்கான ஒரே வழி, இந்த நோக்கத்திற்காக ஒரு தருக்கமல்லாத
குறியை (அதாவது, !!a{function} அல்லது !!{predicate} ஐ) குறிப்பிடுவது.
பொருட்களத்தின் உட்கணங்கள், அதன் மீதான உறவுகள், அல்லது
!!{domain} இலிருந்து !!{domain} க்குச் செல்லும் சார்புகள் அனைத்தும்
ஒரு பண்பைக் கொண்டுள்ளன என்று இரண்டாம்தர அனைத்தளவையிடல் கூற
உதவுகிறது. முதல்தரத் தருக்கத்தில், மொழியின் தருக்கமல்லாத குறிகளில்
ஒன்றுக்கு அளிக்கப்பட்ட உட்கணங்கள், உறவுகள் அல்லது சார்புகள் ஒரு
பண்பைக் கொண்டுள்ளன என்று மட்டுமே நம்மால் கூற முடியும். ஒரு பண்பைக்
கொண்ட உட்கணங்கள், உறவுகள், சார்புகள் உள்ளன என்றோ, அவை அனைத்தும்
அந்தப் பண்பைக் கொண்டுள்ளன என்றோ கூறும்போது, அந்தப் பண்பைக்
குறிப்பிடுவதிலும் இரண்டாம்தர அளவையிடலைப் பயன்படுத்தலாம். இதனால்
முதல்தரத் தருக்கத்தில் வரையறுக்க முடியாத உறவுகளை வரையறுக்கவும்,
முதல்தரத் தருக்கத்தில் வெளிப்படுத்த முடியாத பொருட்களப் பண்புகளை
வெளிப்படுத்தவும் முடிகிறது.
\end{explain}

\begin{defn}
$\Struct{M}$ என்பது மொழி~$\Lang{L}$ க்கான !!a{structure} என்க.
கட்டற்ற மாறிகள் $x$,~$y$ மட்டுமே உள்ள ஏதோ ஒரு
!!{formula}~$!A_R(x, y)$ இருந்து, $s(x) = a$, $s(y) = b$ எனும்போது
$R(a, b)$ பொருந்துவது (அதாவது, $\tuple{a,b} \in R$) அப்போதும்
அப்போது மட்டுமே $\Sat{M}{!A_R(x, y)}[s]$ எனில், ஓர் உறவு
~$R \subseteq \Domain{M}^2$ ஆனது~$\Lang{L}$ இல்
\emph{வரையறுக்கத்தக்கது}.
\end{defn}

\begin{ex}
முதல்தரத் தருக்கத்தில் சமனுறவு~$\Id{\Domain{M}}$ ஐ (அதாவது,
$\Setabs{\tuple{a,a}}{a \in \Domain{M}}$ ஐ) $\eq[x][y]$ என்ற
வாய்பாட்டால் வரையறுக்கலாம். இரண்டாம்தரத் தருக்கத்தில் இந்த உறவை
\emph{~$\eq$ இல்லாமல்} வரையறுக்கலாம். $a$, $b$ ஆகியவை
~$\Domain{M}$ இன் ஒரே !!{element} ஆக இருந்தால், அவை
~$\Domain{M}$ இன் அதே உட்கணங்களின் !!{element}s ஆக இருக்கும்
(கணங்கள் அவற்றின் !!{element}s ஆல் தீர்மானிக்கப்படுவதால்).
மறுதிசையில், $a$, $b$ வேறுபட்டால், அவை அதே உட்கணங்களின்
!!{element}s அல்ல: எடுத்துக்காட்டாக, $a \neq b$ எனில்
$a \in \{a\}$ ஆனால் $b \notin \{a\}$. எனவே “~$\Domain{M}$ இன்
அதே உட்கணங்களின் !!{element}s ஆக இருப்பது” என்பது $a$, $b$
இடையே அப்போதும் அப்போது மட்டுமே $a = b$ ஆகும்போது பொருந்தும் உறவு.
~$\Domain{M}$ இன் எல்லா உட்கணங்கள்மீதும் அளவையிட முடிவதால், அந்த
உறவை இரண்டாம்தரத் தருக்கத்தில் வெளிப்படுத்தலாம். ஆகவே பின்வரும்
!!{formula}, $\Id{\Domain{M}}$ ஐ வரையறுக்கிறது:
\[
\lforall[X][(X(x) \liff X(y))]
\]
\end{ex}

\begin{prob}
$\lforall[X][(X(x) \lif X(y))]$ ஆனது $\Id{\Domain{M}}$ ஐ
வரையறுக்கிறது என்று காட்டுக (கவனிக்க: $\liff$ அன்று,~$\lif$!).
\end{prob}

\begin{ex}
$R$ ஓர் இரண்டிட !!{predicate} எனில், $\Assign{R}{M}$ என்பது
~$\Domain{M}$ மீதான ஓர் இரண்டிட உறவு. சற்றுக் குழப்பமளிக்கக்கூடிய
வகையில்,~$R$ ஐ !!{predicate}~$R$ க்கும், உறவு
~$\Assign{R}{M}$ க்கும் பயன்படுத்துவோம். ~ $R$ இன்
\emph{இடைமாறல் மூடல்}~$R^*$ என்பது, ஏதோ $c_1$, \dots, $c_k$
க்கு $R(a,c_1)$, $R(c_1, c_2)$, \dots, $R(c_k,b)$ பொருந்தினால்,
எனில் மட்டுமே, $a$, $b$ இடையே பொருந்தும் உறவு. இதில் $k = 0$
என்ற நிலையும் அடங்கும்; அதாவது $R(a,b)$ பொருந்தினால்
$R^*(a,b)$ உம் பொருந்தும். இதன் பொருள் $R \subseteq R^*$.
உண்மையில்,~$R$ ஐக் கொண்டதும் இடைமாறுவதுமான மிகச் சிறிய உறவு
~$R^*$. ~ $X$ என்பது~$R$ ஐக் கொண்ட ஓர் இடைமாறும் உறவு என்று
இரண்டாம்தரத் தருக்கத்தில் கூறலாம்:
\begin{multline*}
  !B_R(X) \ident \lforall[x][\lforall[y][(R(x,y) \lif X(x, y))]] \land {}\\
\lforall[x][\lforall[y][\lforall[z][((X(x,y) \land X(y,z)) \lif X(x,
      z))]]].
\end{multline*}
முதல் இணைவு $R \subseteq X$ என்று கூறுகிறது; இரண்டாவது~$X$
இடைமாறுகிறது என்று கூறுகிறது.

~$X$ இவ்வாறான மிகச் சிறிய உறவு என்று கூறுவது,~$R$ ஐக் கொண்டதும்
இடைமாறுவதுமான ஒவ்வோர் உறவிலும் அது தானும் அடங்கும் என்று கூறுவதாகும்.
எனவே~$R$ இன் இடைமாறல் மூடலைப் பின்வரும் !!{formula} மூலம்
வரையறுக்கலாம்:
\[
R^*(X) \ident !B_R(X) \land \lforall[Y][(!B_R(Y) \lif
  \lforall[x][\lforall[y][(X(x, y) \lif Y(x,y))]])].
\]
$\Sat{M}{R^*(X)}[s]$ என்பது $s(X) = R^*$ எனில், எனில் மட்டுமே.
~$R$ இன் இடைமாறல் மூடலை முதல்தரத் தருக்கத்தில் வெளிப்படுத்த
முடியாது.
\end{ex}

\end{document}
